\documentclass[11pt,a4paper]{article}
\usepackage[margin=1in]{geometry}
\usepackage{hyperref}
\usepackage{listings}
\usepackage{xcolor}
\usepackage{graphicx}
\usepackage{enumitem}
\usepackage{fancyhdr}
\usepackage{titlesec}
\usepackage{dirtree}

\hypersetup{
    colorlinks=true,
    linkcolor=blue!70!black,
    urlcolor=blue!70!black
}

\definecolor{codebg}{HTML}{F5F5F5}
\definecolor{codeframe}{HTML}{CCCCCC}

\lstset{
    backgroundcolor=\color{codebg},
    frame=single,
    rulecolor=\color{codeframe},
    basicstyle=\ttfamily\small,
    breaklines=true,
    columns=fullflexible,
    keepspaces=true,
    showstringspaces=false
}

\pagestyle{fancy}
\fancyhf{}
\fancyhead[L]{Windows Migration Tool}
\fancyhead[R]{User Manual}
\fancyfoot[C]{\thepage}

\title{\textbf{Windows Migration Tool}\\[0.5em]\large User Manual}
\author{}
\date{May 2026}

\begin{document}

\maketitle
\thispagestyle{empty}
\vfill
\begin{center}
\textit{A complete guide to exporting and restoring your Windows environment.}
\end{center}
\newpage

\tableofcontents
\newpage

%----------------------------------------------------------------------
\section{Introduction}

The Windows Migration Tool is a PowerShell-based application (with an optional GUI) that helps you move your environment from one Windows 10 or 11 machine to another. It also works as a \textbf{standalone WSL backup and restore tool}---back up your Linux distributions on a schedule or before risky changes, and restore them in minutes.

It automates the tedious parts of PC migration:

\begin{itemize}
    \item \textbf{Backing up and restoring WSL distributions}---full filesystem export as \texttt{.tar} (WSL~1) or \texttt{.vhdx} (WSL~2), plus \texttt{.wslconfig}. Use it as a dedicated WSL backup tool or as part of a full machine migration.
    \item Inventorying all installed software (Win32 and Store apps)
    \item Retrieving license keys (Windows, Office, registry-stored serials)
    \item Exporting winget packages for bulk reinstall
    \item Generating a self-contained restore script
\end{itemize}

%----------------------------------------------------------------------
\section{System Requirements}

\begin{itemize}
    \item Windows 10 (version 1709 or later) or Windows 11 on the \textbf{source} machine
    \item Windows 11 on the \textbf{destination} machine
    \item PowerShell 5.1 or later
    \item \textbf{Administrator privileges} (the tool must be run as Admin)
    \item An external drive or network share with sufficient space for WSL exports
    \item winget (pre-installed on Windows 11; on Windows 10, install from the \href{https://apps.microsoft.com/detail/9NBLGGH4NNS1}{Microsoft Store} or \href{https://github.com/microsoft/winget-cli/releases}{GitHub})
\end{itemize}

%----------------------------------------------------------------------
\section{Installation}

\subsection{Option A: Download from GitHub Releases}

\begin{enumerate}
    \item Go to the \href{https://github.com/rory/PCmigrate/releases}{Releases} page on GitHub.
    \item Download \texttt{PCmigrate\_Setup.exe} (installer) or \texttt{PCmigrate\_Portable.zip} (no install required).
    \item For the installer: run it---it creates a Start Menu shortcut and an optional desktop icon.
    \item For portable: extract the zip and run \texttt{PCmigrate.cmd}.
\end{enumerate}

\noindent\textbf{Note:} Windows SmartScreen may block the unsigned installer. Click \textbf{``More info''} then \textbf{``Run anyway''} to proceed.

\subsection{Option B: Build the Installer Yourself}

\begin{enumerate}
    \item Download and install \href{https://jrsoftware.org/isinfo.php}{Inno Setup 6}.
    \item Compile \texttt{installer.iss} to produce \texttt{PCmigrate\_Setup.exe}.
    \item Run the installer---it creates a Start Menu shortcut and an optional desktop icon.
\end{enumerate}

\subsection{Option C: Portable (No Install, from Source)}

Simply double-click \texttt{PCmigrate.cmd} or run \texttt{PCmigrate-GUI.ps1} directly from an elevated PowerShell prompt. No installation is required.

%----------------------------------------------------------------------
\section{Usage}

\subsection{GUI Mode}

Launch the application from the Start Menu, the desktop shortcut, or by running \texttt{PCmigrate.cmd}.

\begin{enumerate}
    \item Click \textbf{Browse} to select your external drive or output folder.
    \item Click \textbf{Export} to back up this machine (export only).
    \item Click the \textbf{{\footnotesize$\nabla$}} arrow next to Export for options: ``Export Only'', ``Export + Create Restore Bundle'', or ``WSL Only''.
    \item On the destination machine, click \textbf{Browse} to locate the export folder, then click \textbf{Restore}.
    \item Click \textbf{Cancel} at any time to stop a running operation.
\end{enumerate}

The GUI features:
\begin{itemize}
    \item Dark (Catppuccin-style) color scheme
    \item Prominent path label with helper text
    \item Live log output showing progress
    \item Progress bar
    \item Cancel button to abort operations
    \item Automatic self-elevation to Administrator
    \item Non-blocking UI (export/restore runs in background threads)
\end{itemize}

\subsection{Command-Line Mode}

\subsubsection{Exporting (On the Old Machine)}

\begin{lstlisting}[language=bash]
# Export to an external drive
.\Migrate-Machine.ps1 -OutputPath "E:\PCmigrate"

# Export and create a self-contained restore zip
.\Migrate-Machine.ps1 -OutputPath "E:\PCmigrate" -Bundle

# WSL only (skip software inventory and license keys)
.\Migrate-Machine.ps1 -OutputPath "E:\PCmigrate" -WslOnly

# Or use the default (Desktop\PCmigrate)
.\Migrate-Machine.ps1
\end{lstlisting}

\subsubsection{Restoring (On the New Machine)}

Move the external drive to the new machine, then run:

\begin{lstlisting}[language=bash]
E:\PCmigrate\Restore-Machine.ps1
\end{lstlisting}

Or specify the import path manually:

\begin{lstlisting}[language=bash]
.\Restore-Machine.ps1 -ImportPath "E:\PCmigrate"
\end{lstlisting}

%----------------------------------------------------------------------
\section{Migration Workflow}

The recommended end-to-end workflow:

\begin{enumerate}
    \item Plug an external drive into the \textbf{old machine}.
    \item Run the export (GUI or command line) targeting the external drive.
    \item Wait for the export to complete (WSL archives may take time for large distros).
    \item Move the drive to the \textbf{new machine}.
    \item Run the restore script from the export folder.
    \item Complete the post-restore steps (see Section~\ref{sec:postrestore}).
\end{enumerate}

%----------------------------------------------------------------------
\section{Using as a WSL Backup \& Restore Tool}

PCmigrate doubles as a dedicated WSL backup tool. Use the \texttt{-WslOnly} flag (or the ``WSL Only'' option in the GUI) to skip software inventory and license keys and focus exclusively on your Linux distributions.

\subsection{Why Back Up WSL?}

\begin{itemize}
    \item \textbf{Before risky changes} --- experimenting with system packages, kernel upgrades, or distro upgrades? Take a snapshot first.
    \item \textbf{Scheduled backups} --- run the export on a cron-like schedule (Task Scheduler) to maintain rolling backups of your dev environment.
    \item \textbf{Machine migration} --- moving to a new laptop? Export your distros and import them on the other side in minutes.
    \item \textbf{Disaster recovery} --- if a distro becomes unbootable, restore from the last good export.
\end{itemize}

\subsection{WSL Backup Workflow}

\begin{lstlisting}[language=bash]
# Back up all WSL distros to an external drive
.\Migrate-Machine.ps1 -OutputPath "E:\WSL-Backup" -WslOnly

# Restore on the same or a different machine
E:\WSL-Backup\Restore-Machine.ps1
\end{lstlisting}

The export produces \texttt{.tar} files for WSL~1 distros and \texttt{.vhdx} files for WSL~2 distros (on Windows~11 build 22000+). The restore script handles both formats automatically and prints the commands needed to set your default user afterward.

%----------------------------------------------------------------------
\section{Output Structure}

The export produces a self-contained folder with the following layout:

\begin{lstlisting}
PCmigrate/
  license_keys.txt          # Windows/Office/software keys
  installed_software.csv    # Full inventory (machine-readable)
  installed_software.txt    # Full inventory (human-readable)
  installed_software.html   # Interactive checklist with download links
  winget_packages.json      # For 'winget import'
  AppData/
    Roaming_Firefox.zip     # (example) app settings/data
    Local_VSCode.zip        # (example)
  WSL/
    Ubuntu.tar              # (example) WSL 1 distro archive
    Debian.vhdx             # (example) WSL 2 distro disk image
    .wslconfig              # WSL global config
  Restore-Machine.ps1       # Run this on the new machine
  migration_log_*.txt       # Export log
\end{lstlisting}

\subsection{File Descriptions}

\begin{description}[style=nextline]
    \item[\texttt{license\_keys.txt}] Contains product keys retrieved from BIOS/UEFI, Office \texttt{ospp.vbs}, and a registry scan for common key names (Serial, LicenseKey, ProductKey, CDKey, etc.).
    \item[\texttt{installed\_software.csv}] Machine-readable inventory of all installed applications with names, versions, and download URLs.
    \item[\texttt{installed\_software.txt}] Human-readable version of the software inventory.
    \item[\texttt{installed\_software.html}] Interactive checklist with checkboxes, download/search links, and license keys (hidden by default with a copy button). Open in any browser on the new machine.
    \item[\texttt{winget\_packages.json}] Package list for \texttt{winget import} to automate bulk reinstallation.
    \item[\texttt{AppData/*.zip}] Zip archives of application settings and data from \texttt{\%APPDATA\%} and \texttt{\%LOCALAPPDATA\%}. Named \texttt{Roaming\_AppName.zip} or \texttt{Local\_AppName.zip}. Only includes folders matching installed apps, up to 500~MB each.
    \item[\texttt{WSL/*.tar}] Each WSL 1 distribution exported as a tar archive.
    \item[\texttt{WSL/*.vhdx}] Each WSL 2 distribution exported as a virtual disk image (Windows 11 only).
    \item[\texttt{WSL/.wslconfig}] Global WSL configuration file.
    \item[\texttt{Restore-Machine.ps1}] Auto-generated restore script that uses \texttt{\$PSScriptRoot} to locate sibling data files.
    \item[\texttt{migration\_log\_*.txt}] Timestamped log of the export process.
\end{description}

%----------------------------------------------------------------------
\section{What the Restore Script Does}

The generated \texttt{Restore-Machine.ps1} script performs:

\begin{enumerate}
    \item \textbf{winget import} --- Reads \texttt{winget\_packages.json} and installs all listed packages.
    \item \textbf{WSL import} --- Runs \texttt{wsl --import} for each \texttt{.tar} or \texttt{.vhdx} archive in the \texttt{WSL/} folder, restoring your Linux distributions.
    \item \textbf{App data restore} --- Extracts each zip in \texttt{AppData/} back to the corresponding \texttt{\%APPDATA\%} or \texttt{\%LOCALAPPDATA\%} folder (skips folders that already exist).
\end{enumerate}

After completion, the script prints a reminder showing the exact commands needed to set your default user for each restored WSL distro.

The script is self-contained: it uses \texttt{\$PSScriptRoot} so it works from wherever the export folder is located on the new machine.

%----------------------------------------------------------------------
\section{Post-Restore Steps}
\label{sec:postrestore}

After running the restore script, it prints a reminder with the exact commands you need. Complete these steps:

\begin{enumerate}
    \item \textbf{Set WSL default users:} Imported distros default to root. The restore script prints the commands, but for reference:\\
    \texttt{<distro> config --default-user <username>}\\
    For example: \texttt{Ubuntu config --default-user yourname}
    \item \textbf{Install remaining apps:} Check \texttt{installed\_software.csv} for applications not available via winget and install them manually.
    \item \textbf{Re-enter license keys:} Open \texttt{license\_keys.txt} and activate software that requires manual key entry.
    \item \textbf{Restore app-specific settings:} Any application configurations not covered by this tool (browser profiles, IDE settings, etc.) must be transferred separately.
\end{enumerate}

%----------------------------------------------------------------------
\section{Limitations}

\begin{itemize}
    \item License keys are only retrievable if stored in BIOS/UEFI or the Windows registry. Digital licenses tied to a Microsoft account transfer automatically and will not appear in the export.
    \item Some winget packages may fail to import if the source or package ID has changed since the export.
    \item WSL exports include the full filesystem---large distributions will produce large \texttt{.tar} files. Ensure your external drive has sufficient free space.
    \item Framework and runtime packages from the Microsoft Store are filtered out to reduce noise, but some may still appear.
\end{itemize}

%----------------------------------------------------------------------
\section{Troubleshooting}

\begin{description}[style=nextline]
    \item[``Access Denied'' or permission errors]
    Ensure you are running the tool as Administrator. The GUI self-elevates automatically; for command-line use, right-click PowerShell and select ``Run as administrator.''

    \item[WSL export fails]
    Verify that WSL is installed and that your distributions are in a stopped state. Run \texttt{wsl --shutdown} before exporting.

    \item[winget import partially fails]
    Some packages may have been removed or renamed. Check the log file and install those packages manually.

    \item[Export folder is very large]
    WSL distributions contain full Linux filesystems. Consider cleaning up unnecessary files inside WSL before exporting (e.g., package caches, temp files).
\end{description}

%----------------------------------------------------------------------
\section{Continuous Integration \& Releases}

The project includes GitHub Actions workflows for automated builds:

\begin{description}[style=nextline]
    \item[\texttt{.github/workflows/build-installer.yml}] Runs on every push and pull request to \texttt{main}. Compiles the Inno Setup installer and uploads it as a workflow artifact.
    \item[\texttt{.github/workflows/release.yml}] Triggered when a version tag (e.g., \texttt{v1.0.0}) is pushed. Builds the installer, packages a portable zip, and publishes both as GitHub Release assets.
\end{description}

To create a release:

\begin{lstlisting}[language=bash]
git tag v1.0.0
git push origin v1.0.0
\end{lstlisting}

This produces a GitHub Release with two downloadable assets:
\begin{itemize}
    \item \texttt{PCmigrate\_Setup.exe} --- full Windows installer
    \item \texttt{PCmigrate\_Portable.zip} --- portable scripts (no installation required)
\end{itemize}

%----------------------------------------------------------------------
\section{Project Files Reference}

\begin{description}[style=nextline]
    \item[\texttt{Migrate-Machine.ps1}] Main export script. Accepts \texttt{-OutputPath} parameter.
    \item[\texttt{PCmigrate-GUI.ps1}] WPF-based graphical interface.
    \item[\texttt{PCmigrate.cmd}] CMD launcher that starts the GUI without a console flash.
    \item[\texttt{installer.iss}] Inno Setup 6 script for building the Windows installer.
    \item[\texttt{icon.ico}] Application icon used by the installer.
    \item[\texttt{.github/workflows/build-installer.yml}] CI workflow for building the installer on push/PR.
    \item[\texttt{.github/workflows/release.yml}] Release workflow for publishing tagged versions.
\end{description}

\end{document}
